\section*{ID6656: They Both Come from Fourier-Like Radial Symmetry}
\paragraph{Metadata.}
PiID: 6643193950979 | RTEID: RTE:P25846.7261:T4.8883 | UUID: 75e842b9-f838-5513-b536-0de924afa0f3

\paragraph{Canonical Statement.}
\textbackslash{}[ I(\textbackslash{}theta) \textbackslash{}sim \textbackslash{}left| \textbackslash{}sum\_\{n=0\}\textasciicircum{}\{N-1\} e\textasciicircum{}\{i k r \textbackslash{}cos(\textbackslash{}theta - \textbackslash{}theta\_n)\} \textbackslash{}right|\textasciicircum{}2 \textbackslash{}]

\paragraph{Canonical Equation.}
\[
I(\theta) \sim \left| \sum_{n=0}^{N-1} e^{i k r \cos(\theta - \theta_n)} \right|^2
\]

\paragraph{Dictionary.}
They Both Come from Fourier-Like Radial Symmetry — I(θ) ∼ | ∑\_n=0\textasciicircum{}N-1 e\textasciicircum{}i k r cos(θ - θ\_n) |\textasciicircum{}2

\paragraph{Encyclopedia.}
They Both Come from Fourier-Like Radial Symmetry is a series / summation, written I(θ) ∼ | ∑\_n=0\textasciicircum{}N-1 e\textasciicircum{}i k r cos(θ - θ\_n) |\textasciicircum{}2. It is expressed as a sum or series over modes or indices, superposing discrete contributions. It is built from θ, I, n, N, k, r. Within the Trawin Zero Point registry it serves as one element of the framework's mathematical narrative.
